\section*{Bab \ref{ch:g4:measure} --- Pengukuran dan Keliling}

\begin{solution}{exo:g4:measure:1}
$7$ m $= 700$ cm; \quad $3.2$ km $= 3\,200$ m; \quad
$85$ mm $= 8.5$ cm; \quad $4\,500$ m $= 4.5$ km.
\end{solution}

\begin{solution}{exo:g4:measure:2}
$2$ kg $= 2\,000$ g; \quad $1\,800$ g $= 1$ kg $800$ g; \quad
$3.5$ L $= 350$ cL; \quad $25$ cL $= 0.25$ L.
\end{solution}

\begin{solution}{exo:g4:measure:3}
Persegi panjang: $2 \times (9 + 5) = 28$ cm. Persegi:
$4 \times 7.5 = 30$ cm. Segitiga: $6 + 8 + 10 = 24$ cm.
\end{solution}

\begin{solution}{exo:g4:measure:4}
Panjang $+$ lebar $= 26 \div 2 = 13$ cm, jadi lebarnya
$13 - 8 = 5$ cm.
\end{solution}

\begin{solution}{exo:g4:measure:5}
Banyak bentuk L yang berhasil; untuk L seluas $10$ petak yang
tersusun dari persegi panjang $4 \times 2$ ditambah kaki
$2 \times 1$, menelusuri tepinya memberi $16$ satuan. (Setiap gambar
yang benar dan hitungan yang teliti itu benar; \omterm{def:g4:measure:perimeter}{kelilingnya} bergantung pada bentuk L yang dipilih.)
\end{solution}

\begin{solution}{exo:g4:measure:6}
$3 \times 4$: \omterm{def:g4:measure:perimeter}{keliling} $14$. $2 \times 6$: \omterm{def:g4:measure:perimeter}{keliling} $16$.
($1 \times 12$: \omterm{def:g4:measure:perimeter}{keliling} $26$.) Makin dekat ke bentuk persegi, makin
kecil \omterm{def:g4:measure:perimeter}{kelilingnya} untuk \omterm{def:g4:measure:area}{luas} yang sama.
\end{solution}

\begin{solution}{exo:g4:measure:7}
Setengah dari persegi panjang $6 \times 3$: $18 \div 2 = 9$ petak.
\end{solution}

\begin{solution}{exo:g4:measure:8}
$10{:}20 \to 12{:}45$: $2$ jam $25$ menit. $7{:}35 \to 9{:}10$:
$1$ jam $35$ menit. Film: $20{:}30 + 1$ jam $50 = 22{:}20$.
\end{solution}

\begin{solution}{exo:g4:measure:9}
Satu putaran: $2 \times (60 + 45) = 210$ m. Tiga putaran:
$3 \times 210 = 630$ m.
\end{solution}

\begin{solution}{exo:g4:measure:10}
\omterm{def:g4:measure:perimeter}{Keliling}: $4 \times 85 = 340$ m; dikurangi gerbangnya:
$340 - 4 = 336$ m pagar.
\end{solution}

\begin{solution}{exo:g4:measure:11}
Kedua pernyataan itu \emph{salah}. \omterm{def:g4:measure:perimeter}{Keliling} sama, \omterm{def:g4:measure:area}{luas} berbeda:
persegi panjang $4 \times 2$ dan bangun L pada gambar bab ini
(keduanya berkeliling $12$; luasnya $8$ dan $6$). \omterm{def:g4:measure:area}{Luas} sama,
\omterm{def:g4:measure:perimeter}{keliling} berbeda: persegi $4 \times 4$ dan persegi panjang
$8 \times 2$ (keduanya berluas $16$; \omterm{def:g4:measure:perimeter}{kelilingnya} $16$ dan $20$).

\end{solution}
